\lecture{12}{Relational Deep Learning}{Thursday, October 29, 2026}\label{ch:lec12}

% Local notation (lecture 12)
\providecommand{\vt}{\mathbf{t}}

% Local TikZ styles (prefix lec12)
\tikzset{
  lec12usr/.style={circle, draw=black!70, thick, fill=Orchid!35, minimum size=5mm,
                   inner sep=0pt, font=\scriptsize},
  lec12sale/.style={rectangle, rounded corners=1.5pt, draw=black!70, thick,
                    fill=Goldenrod!35, inner sep=1.5pt, font=\tiny},
  lec12prod/.style={circle, draw=black!70, thick, fill=SkyBlue!45, minimum size=5mm,
                    inner sep=0pt},
  lec12churn/.style={rectangle, rounded corners=1.5pt, draw=RoyalBlue!70, fill=RoyalBlue!15,
                     font=\tiny\sffamily, minimum width=9mm, inner sep=1.5pt},
  lec12act/.style={rectangle, rounded corners=1.5pt, draw=Goldenrod!90!black,
                   fill=Goldenrod!20, font=\tiny\sffamily, minimum width=9mm, inner sep=1.5pt},
  lec12time/.style={font=\tiny\sffamily, black!60},
  lec12att/.style={densely dashed, black!60, -{Stealth[length=1.5mm]}},
  lec12tr/.style={circle, draw=black, fill=Cyan!45, minimum size=5.2mm, inner sep=0pt,
                  font=\scriptsize},
  lec12pr/.style={circle, draw=black, fill=Orange!55, minimum size=5.2mm, inner sep=0pt,
                  font=\scriptsize},
  lec12cu/.style={circle, draw=black, fill=ForestGreen!45, minimum size=5.2mm,
                  inner sep=0pt, font=\scriptsize},
  lec12oe/.style={thick, Orange!85!black, {Stealth[length=1.4mm]}-{Stealth[length=1.4mm]}},
  lec12ge/.style={thick, ForestGreen!80!black, {Stealth[length=1.4mm]}-{Stealth[length=1.4mm]}},
  lec12cap/.style={font=\scriptsize\sffamily, align=center},
  lec12smp/.style={circle, draw=Purple!70, thick, fill=Purple!8, minimum size=4mm,
                   inner sep=0pt},
}

\textbf{Outline.} Relational databases (motivation, why not LLMs) $\to$ generalizing deep
learning to databases (RDL) $\to$ a database is a graph $\to$ relational database graph
(schema graph, relational entity graph) $\to$ predictive tasks (training tables) $\to$
relational GNN (temporal, SQL joins vs.\ edges, benefits) $\to$ RelBench $\to$ GNN vs.\
expert data scientist $\to$ results.

% ---------------------------------------------------------------------
\section{Relational databases}\label{sec:lec12-rdb}

\subsection{Deep learning revolution}\label{sec:lec12-revolution}
\slides{4--6}

\begin{itemize}
  \item Deep learning revolution on \textbf{sequences} (text, e.g.\ the Wikipedia article
    ``Statue of Liberty''; protein chains A and B as amino-acid sequences) and
    \textbf{images} (photos; a Tesla Autopilot tweet).
  \item These breakthroughs are fueled by \textbf{Transformers}: encoder--decoder
    architecture (input/output embedding + positional encoding; $N\times$ multi-head
    attention, add \& norm, feed forward; masked multi-head attention; linear; softmax
    $\to$ output probabilities); \emph{ChatGPT's exam results} (percentile rank jumps
    63rd$\to$99th, 25th$\to$80th, 40th$\to$88th, 22nd$\to$71st, 10th$\to$90th,
    30th$\to$66th).
  \item \dots\ and fueled by \textbf{Data}: The Economist (May 6th--12th 2017), ``The
    world's most valuable resource -- Data and the new rules of competition''.
\end{itemize}

\subsection{Data in relational databases and predictive tasks}\label{sec:lec12-tasks}
\slides{7--10}

\begin{itemize}
  \item \textbf{Data stored in relational databases} (commerce and finance schemas as in
    \cref{ch:lec01}); new: \textbf{Social Media} -- Users, UserSettings, UserProfiles,
    Locations, ExternalAccounts, FollowingRelationships, Posts, PostComments, PostStats,
    PostFavorites, PostTags, Tags, PostCategories, Categories.
  \item \textbf{Predictions on relational data:}
    \begin{itemize}
      \item \textbf{Commerce} (Users(ID) -- Sales(ID, Product\_ID, Session\_ID, User\_ID,
        Price) -- Products(Product\_ID)): which \textbf{products will a user purchase} in
        the next 7 days? Will an \textbf{active user churn} in the next 90 days? What will
        be the \textbf{total sales for each product} in the next 30 days?
      \item \textbf{Finance} (loan, order, trans, card, disp, account, client, district):
        what will be the \textbf{loan outcome} (good vs bad)? What will be the
        \textbf{total number of failed loans} in the next year?
      \item \textbf{Health} (patients 46,520 rows; admissions 58,976 rows; noteevents
        2,083,180 rows; procedures\_icd 240,095 rows; diagnoses\_icd 651,047 rows; linked
        by subject\_id, hadm\_id): will a \textbf{patient return} if discharged from the
        hospital? Which hospital admissions have \textbf{greatest risk to life} in the
        next 24 hours?
    \end{itemize}
\end{itemize}

\subsection{Doing AI is slow \& complex}\label{sec:lec12-slow}
\slides{11--13}

\begin{itemize}
  \item \textbf{Pipeline:} input data cleaning, curation $\to$ target label engineering
    $\to$ feature engineering $\to$ architecture \& hyperparameter search $\to$ ML Ops:
    \textbf{6-12 months of time for a team of data scientists, data engineers and product
    engineers}.
  \item \term{Impedance mismatch}. \textbf{Goal:} learn a \textbf{user churn} model based
    on their sales, purchased products and browsing behavior. Tables Users(ID),
    Sessions(ID, User\_ID), Sales(ID, Product\_ID, Session\_ID, User\_ID, Price),
    Products(Product\_ID), Views(ID, Product\_ID, Session\_ID) are flattened into
    \textbf{one table}: rows = examples, columns = features (\#monthly sales, sum weekly
    item price, \#orders relative to Sunday) + Label. Problems:
    \begin{itemize}
      \item features are chosen \textbf{arbitrarily} (e.g.\ aggregations, time windows);
      \item only a \textbf{limited set of data} is used;
      \item issues with \textbf{point-in-time correctness/information leakage}.
    \end{itemize}
\end{itemize}

\begin{keyidea}[Key question]
  How to leverage the data without going through the longest duration tasks
  (extraction, transformation, loading)? \textbf{We want ML algorithms that can process
  data in its natural form!} (slides~13, 18)
\end{keyidea}

\subsection{Why not use an LLM?}\label{sec:lec12-llm}
\slides{14--17}

\begin{itemize}
  \item \textbf{Some have tried:} Wydmuch, Borchmann, Grali\'nski (Snowflake AI Research),
    \emph{Tackling prediction tasks in relational databases with LLMs}. Process of
    constructing a single example for LLM-based inference: schema (Task Table: Target,
    Timestamp, Customer ID; Customers; Products; Transactions) $\to$ collect the rows
    linked to one task row $\to$ serialize as JSON (\texttt{"Customer ID": 13},
    \dots, \texttt{"Transactions": [\dots]}, \texttt{"Target": 1}).
  \item \textbf{Prompt} (slide~16): one such serialized example (product attributes,
    transactions, prices, timestamps) already fills a whole slide.
  \item \textbf{Why not LLMs?}
    \begin{itemize}
      \item LLMs are not trained to understand and effectively learn from the
        \textbf{relational nature} of Databases;
      \item Database prediction is \textbf{NOT} sequence modeling;
      \item LLMs are \textbf{NOT trained to predict the future};
      \item \textbf{Context size} is \textbf{limited}.
    \end{itemize}
\end{itemize}

\subsection{Modern ML}\label{sec:lec12-modern}
\slides{18--19}

\begin{itemize}
  \item Slide~18 repeats the key question above.
  \item Today we want to design deep learning models that operate on \textbf{relational
    tables}; but the modern deep learning toolbox is designed for \textbf{different type
    of inputs} (text, images, networks).
\end{itemize}

% ---------------------------------------------------------------------
\section{Generalizing deep learning to databases}\label{sec:lec12-generalize}

\subsection{Relational deep learning (RDL)}\label{sec:lec12-rdl}
\slides{21--23}

\begin{itemize}
  \item \textbf{Deep learning on relational tables:} relational tables $\to$
    \term{Relational Deep Learning} $\to$ \textbf{Prediction}.
  \item \textbf{What is RDL?}
    \begin{itemize}
      \item End-to-end \textbf{deep learning on relational tables} (i.e., databases);
      \item works directly on relational tables, no transformations, no feature
        engineering;
      \item casts predictive tasks as \textbf{graph representation learning} problems.
    \end{itemize}
  \item \textbf{Impact and consequences:}
    \begin{itemize}
      \item \textbf{More accurate models} (no feature engineering);
      \item \textbf{More robust models} (model-learned features automatically ``update''
        over time);
      \item \textbf{Shorter time to models} (no mundane ETL work);
      \item \textbf{Simpler infrastructure} (no pipelines, no feature stores, etc.).
    \end{itemize}
\end{itemize}

\subsection{Related work}\label{sec:lec12-related}
\slides{24--25}

\begin{itemize}
  \item \textbf{Tabular ML} (e.g.\ decision trees: root/decision/leaf nodes, sub-tree):
    \begin{itemize}
      \item great for building models on \textbf{one table} -- but most data is not a
        single table!
      \item deep learning is not dominant. \textbf{Hypothesis:} because single table
        already contains features engineered from other tables (\textbf{loss of
        information}).
      \item \textbf{RDL accounts for relations between multiple tables, unlike tabular
        ML!}
    \end{itemize}
  \item \textbf{Statistical Relational Learning} [Friedman et al., 1999]: learns a
    \textbf{distribution (graphical model) over the relational structure} (example:
    Movie(movie-id, title, process, decade, genre), Appears(appears-id, movie, actor,
    roletype), Actor(actor-id, gender); dotted = relationship, arrow = probabilistic
    dependency).
    \begin{itemize}
      \item predictions via world model + inductive logic programming;
      \item only works with categorical variables;
      \item does not scale to large data.
    \end{itemize}
  \item RDL: a \textbf{scalable, expressive inheritor} of Statistical Relational Learning.
\end{itemize}

% ---------------------------------------------------------------------
\section{Relational deep learning: a database is a graph}\label{sec:lec12-dbgraph}
\slides{27--30}

\begin{itemize}
  \item \textbf{Insight: a database is a graph!} Schema Users(ID) -- Sales(ID,
    Product\_ID, Session\_ID, User\_ID, Price) -- Products(Product\_ID) $\Rightarrow$
    graph Users -- Sales (\$20, \$80, \$15, \$20, \$80, \$80, \$15) -- Products (shirt,
    shoe, sunglasses) -- the same figure as in \cref{ch:lec01}.
  \item \textbf{Just do ML on a graph!} ML in the language of graphs:
    \begin{itemize}
      \item \textbf{node-level:} churn, life-time value, next best action;
      \item \textbf{link-level:} product affinity, recommendations;
      \item \textbf{graph-level:} fraud, money laundering.
    \end{itemize}
  \item \textbf{Graph ML problem solving pipeline:} Relational DB $\to$ Graph Problem
    $\to$ Graph ML (neural networks on computation graphs) $\to$ Solutions (predicted
    labels attached to the users).
\end{itemize}

% ---------------------------------------------------------------------
\section{Relational database graph}\label{sec:lec12-rdbgraph}

\subsection{Databases are relational}\label{sec:lec12-relational}
\slides{32--34}

\begin{itemize}
  \item \textbf{Real-world data are stored in databases:} a transaction (User $\to$
    Product) becomes a row of the \textbf{Transaction Table} (Trans ID, User ID, Product
    ID, Value, Timestamp), e.g.\ (100, 51, 345, \$1000, 1710281001).
  \item \textbf{Databases are often relational:} User ID 51 points to the \textbf{User
    Table} row (51, Male) (columns User ID, Gender); Product ID 345 points to the
    \textbf{Product Table} row (345, Bike) (columns Product ID, Category)
    $\Rightarrow$ \textbf{relational databases as heterogeneous graphs}
    (\cref{ch:lec09}).
  \item \textbf{Mathematically:} a database is a set of tables $\calT$ and links between
    tables $\mathcal{L}$:
\end{itemize}
\begin{equation}\label{eq:lec12-db}
  \calT=\{T_1,\dots,T_n\},\qquad \mathcal{L}\subseteq\calT\times\calT
\end{equation}

\subsection{Schema graph}\label{sec:lec12-schema}
\slides{35--36}

\begin{itemize}
  \item The \term{schema graph} represents the \textbf{high-level structure} of the
    heterogeneous graph. Example: Products(ProductID, Description, Image, Size),
    Transactions(TransactionID, ProductID, \textbf{Timestamp}, CustomerID, Price),
    Customers(CustomerID, Name); edges ``Bought which'' (Products--Transactions) and
    ``Bought by'' (Transactions--Customers). Second example: orders, orders\_items,
    customers, items.
  \item \raggedright The schema graphs can be \textbf{more complex}: \textbf{rel-trial} schema --
    studies (key nct\_id), outcomes, outcome\_analyses, drop\_withdrawals, designs,
    eligibilities, reported\_event\_totals, interventions, interventions\_studies,
    facilities, facilities\_studies, sponsors, sponsors\_studies, conditions,
    conditions\_studies (columns typed numerical / categorical / text / timestamp).
\end{itemize}

\subsection{Relational entity graph}\label{sec:lec12-entity}
\slides{37--39}

\begin{itemize}
  \item \term{Relational entity graph}: create connections via \textbf{primary-foreign
    keys} (the Users--Sales--Products graph above / \cref{ch:lec01}).
\end{itemize}

\begin{definition}[Relational entity graph]\label{def:lec12-reg}
  Given a database of tables $\calT=\{T_1,\dots,T_n\}$ and relations
  $\mathcal{L}\subseteq\calT\times\calT$, where each table is a set of entities $v\in T_i$
  possessing a \textbf{primary key} and optional \textbf{foreign keys}, the relational
  entity graph is such that
  \begin{itemize}
    \item the \textbf{set of nodes} is defined by all rows in all tables:
      \begin{equation}\label{eq:lec12-nodes}
        \calV=\bigcup_{T\in\calT}T
      \end{equation}
    \item the \textbf{set of edges} is defined by connecting two entities $v_1,v_2$ whose
      primary and foreign keys match.
  \end{itemize}
\end{definition}

\begin{itemize}
  \item \textbf{Entities also have features} (different than KGs, \cref{ch:lec10}): e.g.\
    Gender (Male), Category (Bike), Value (\$1000), Timestamp.
\end{itemize}

% ---------------------------------------------------------------------
\section{Predictive tasks in relational databases}\label{sec:lec12-predtasks}

\subsection{Overview and defining a task}\label{sec:lec12-define}
\slides{41--42}

\begin{itemize}
  \item \textbf{Overview of the approach:} (a) \textbf{define tasks} (``How much will each
    customer buy in the next 90 days?'') $\to$ (b) \textbf{training table generation}:
    gather training targets from historical data (Training Table: 90-day Target,
    Customer\_ID, SeedTime) $\to$ (c) \textbf{link to relational tables} (Training
    Table.User\_ID -- Customers.Customer\_ID -- Transactions.Customer\_ID,
    Transactions.Product\_ID -- Products.Product\_ID).
    \textbf{Next:} define task(s); relational deep learning.
  \item \textbf{Example:} predict whether a user is going to churn in the next 30 days?
  \item \textbf{Most tasks are temporal:}
    \begin{itemize}
      \item customer's \textbf{label changes over time};
      \item database changes over time -- for every time, a feature needs to be
        recomputed.
    \end{itemize}
  \item Temporal tasks are especially challenging because features are time dependent:
    for every time, a feature needs to be recomputed; entity's label can change between
    time steps.
  \item \textbf{To define a task, we need:} \textbf{Entity}, \textbf{Label},
    \textbf{Time}.
\end{itemize}

\subsection{Training table}\label{sec:lec12-training}
\slides{43--45}

\begin{definition}[Training table]\label{def:lec12-train}
  \term{Training Table}: a special table containing training labels, with rows
  \textbf{(entity ID, time, labels)}; tasks: classification, regression, multi-class.
\end{definition}

\begin{itemize}
  \item The \textbf{time column} is essential for \textbf{temporal prediction tasks}:
    \begin{itemize}
      \item an entity may have different labels at different times;
      \item only use information up to the time of label.
    \end{itemize}
  \item \textbf{Example} (slide~44), rows (Entity ID, Timestamp, Label): (99, 10172024, 1),
    (99, 10182024, 1), \dots, (100, 10172024, 1), (100, 10182024, 0), \dots\ -- entity
    100 changes label between the two times.
  \item \textbf{Example -- churn:} schema Users--Sales--Products; task: predict whether a
    user is going to churn = \textbf{zero sales in the next 30 days}; training table
    \textbf{(user, time, churn label)}.
\end{itemize}

\subsection{RDL pipeline and connecting the training table}\label{sec:lec12-pipeline}
\slides{46--48}

\begin{itemize}
  \item \textbf{Relational Deep Learning pipeline} (slide~46, repeated on slide~50):
    \begin{itemize}
      \item (a) rel.\ tables with training table (Customer LTV: 90-day LTV, User\_ID,
        SeedTime; Customers, Transactions, Products);
      \item (b) entities (rows) linked by foreign keys;
      \item (c) relational entity graph: products 1, 2, 3 -- transactions -- customers --
        training-table nodes (A, B);
      \item (d) graph neural network: computation graph rooted at each training-table node
        (B, A), messages flow from products through transactions and customers.
    \end{itemize}
  \item Relational entity graph (slide~47 = slide~37).
  \item \textbf{Connect the training table:} \textbf{training labels} together with
    \textbf{timestamps} are \textbf{attached} to the graph -- each user gets one label
    node per time (\cref{fig:lec12-labels}a).
\end{itemize}

\begin{figure}[htbp]
  \centering
  \begin{tikzpicture}[scale=0.85, transform shape]
    % (a) training table attached to the entity graph
    \begin{scope}
      \foreach \i/\y/\t/\l/\s in {1/3.4/May '23/Churn/lec12churn,
                                  2/2.95/June '23/Active/lec12act,
                                  3/2.4/May '23/Churn/lec12churn,
                                  4/1.95/June '23/Active/lec12act,
                                  5/1.35/June '23/Active/lec12act,
                                  6/0.8/April '23/Active/lec12act,
                                  7/0.35/May '23/Churn/lec12churn,
                                  8/-0.1/June '23/Active/lec12act} {
        \node[\s] (l\i) at (0.9,\y) {\l};
        \node[lec12time, anchor=east] at (0.38,\y) {\t};
      }
      \node[lec12usr] (u1) at (2.6,3.175) {$u_1$};
      \node[lec12usr] (u2) at (2.6,2.175) {$u_2$};
      \node[lec12usr] (u3) at (2.6,1.35) {$u_3$};
      \node[lec12usr] (u4) at (2.6,0.35) {$u_4$};
      \draw[lec12att] (u1) -- (l1.east); \draw[lec12att] (u1) -- (l2.east);
      \draw[lec12att] (u2) -- (l3.east); \draw[lec12att] (u2) -- (l4.east);
      \draw[lec12att] (u3) -- (l5.east);
      \draw[lec12att] (u4) -- (l6.east); \draw[lec12att] (u4) -- (l7.east);
      \draw[lec12att] (u4) -- (l8.east);
      \foreach \i/\y/\p in {1/3.4/20, 2/2.8/80, 3/2.2/15, 4/1.6/20, 5/1.0/80, 6/0.4/80, 7/-0.2/15}
        \node[lec12sale] (s\i) at (4.2,\y) {\$\p};
      \node[lec12prod, label={[font=\scriptsize]right:shirt}] (sh) at (5.8,2.8) {};
      \node[lec12prod, label={[font=\scriptsize]right:shoe}] (so) at (5.8,1.6) {};
      \node[lec12prod, label={[font=\scriptsize]right:sunglasses}] (sg) at (5.8,0.4) {};
      \draw[edge] (u1)--(s1) (u1)--(s2) (u2)--(s3) (u3)--(s4) (u3)--(s5) (u4)--(s6) (u4)--(s7);
      \draw[edge] (s1)--(sh) (s4)--(sh) (s2)--(so) (s5)--(so) (s6)--(so) (s3)--(sg) (s7)--(sg);
      \node[font=\scriptsize] at (2.6,-0.65) {Users};
      \node[font=\scriptsize] at (4.2,-0.65) {Sales};
      \node[font=\scriptsize] at (5.8,-0.65) {Products};
      \node[lec12cap] at (3.0,-1.15) {(a) labels + timestamps attached to the graph (sl.~48)};
    \end{scope}
    % (b) computation graphs of u1
    \begin{scope}[xshift=8.6cm]
      \node[lec12time, anchor=east] at (-0.05,2.8) {May '23};
      \node[lec12churn] (c1) at (0.5,2.8) {Churn};
      \node[lec12usr] (a1) at (1.7,2.8) {$u_1$};
      \node[lec12sale] (a2) at (3.0,2.8) {\$20};
      \node[lec12prod, label={[font=\scriptsize]right:shirt}] (a3) at (4.3,2.8) {};
      \draw[lec12att] (a1) -- (c1);
      \draw[edge] (a1)--(a2)--(a3);
      \node[lec12time, anchor=east] at (-0.05,1.0) {June '23};
      \node[lec12act] (c2) at (0.5,1.0) {Active};
      \node[lec12usr] (b1) at (1.7,1.0) {$u_1$};
      \node[lec12sale] (b2) at (3.0,1.5) {\$20};
      \node[lec12sale] (b3) at (3.0,0.5) {\$80};
      \node[lec12prod, label={[font=\scriptsize]right:shirt}] (b4) at (4.3,1.5) {};
      \node[lec12prod, label={[font=\scriptsize]right:shoe}] (b5) at (4.3,0.5) {};
      \draw[lec12att] (b1) -- (c2);
      \draw[edge] (b1)--(b2)--(b4) (b1)--(b3)--(b5);
      \node[lec12cap] at (2.3,-1.15) {(b) GNN computation graphs of $u_1$ (sl.~51)};
    \end{scope}
  \end{tikzpicture}
  \caption{Training labels with timestamps attached to the relational entity graph, and the
  time-dependent computation graphs of the first user (slides~48, 51).}
  \label{fig:lec12-labels}
\end{figure}

% ---------------------------------------------------------------------
\section{Relational GNN}\label{sec:lec12-rgnn}

\subsection{GNN on the entity graph}\label{sec:lec12-gnn}
\slides{50--52}

\begin{itemize}
  \item Slide~50 repeats the RDL pipeline (slide~46).
  \item \textbf{Node's neighborhood defines a computation graph} (\cref{ch:lec03}). Nodes
    learn how to \emph{optimally} use information from neighbors to obtain enhanced node
    representations (every user, sale and product node carries an embedding).
    \textbf{Example} (\cref{fig:lec12-labels}b): the first user's May '23 (Churn)
    computation graph contains only the \$20 sale $\to$ shirt; the June '23 (Active) one
    also contains the \$80 sale $\to$ shoe.
  \item \textbf{GNNs on temporal graphs} (\cref{fig:lec12-temporal}): (a) schema graph
    (Products, Transactions, Customers), (b) relational entity graph with transactions
    $v_1,\dots,v_6$ ordered in time, (c) computation graphs of customer $v_{11}$ for
    different time $t$.
    \begin{itemize}
      \item The computation graph for each node is \textbf{time-dependent};
      \item Message+Aggregation becomes \textbf{time-dependent};
      \item Sampling over neighbors is \textbf{time-dependent}.
    \end{itemize}
\end{itemize}

\begin{figure}[htbp]
  \centering
  \begin{tikzpicture}[scale=0.85, transform shape]
    % (b) relational entity graph
    \begin{scope}
      \draw[-{Stealth[length=1.5mm]}, black!70] (-0.75,3.3) -- (-0.75,-0.2);
      \node[font=\scriptsize] at (-1.0,3.3) {0};
      \node[font=\scriptsize] at (-1.0,-0.2) {6};
      \node[font=\scriptsize, rotate=90] at (-1.05,1.55) {Time};
      \foreach \i in {1,...,6} \node[lec12tr] (v\i) at (0,{3.3-(\i-1)*0.68}) {$v_{\i}$};
      \foreach \i in {7,8,9} \node[lec12pr] (v\i) at (2.6,{3.3-(\i-7)*0.62}) {$v_{\i}$};
      \foreach \i in {10,11,12} \node[lec12cu] (v\i) at (2.6,{1.2-(\i-10)*0.68}) {$v_{\i}$};
      \draw[lec12oe] (v1)--(v7); \draw[lec12oe] (v2)--(v7); \draw[lec12oe] (v3)--(v8);
      \draw[lec12oe] (v4)--(v8); \draw[lec12oe] (v6)--(v8); \draw[lec12oe] (v5)--(v9);
      \draw[lec12ge] (v1)--(v11); \draw[lec12ge] (v2)--(v10); \draw[lec12ge] (v3)--(v11);
      \draw[lec12ge] (v4)--(v12); \draw[lec12ge] (v5)--(v11); \draw[lec12ge] (v6)--(v12);
      \node[lec12cap] at (1.1,-0.85) {(b) relational entity graph};
    \end{scope}
    % (c) computation graphs for t = 1, 3, 5
    \begin{scope}[xshift=4.6cm]
      \node[font=\small] at (0,3.75) {$t=1$};
      \node[lec12cu] (a11) at (0,3.2) {$v_{11}$};
      \node[lec12tr] (a1) at (0,1.9) {$v_1$};
      \node[lec12pr] (a7) at (0,0.6) {$v_7$};
      \draw[dedge] (a11)--(a1); \draw[dedge] (a1)--(a7);
      \node[font=\small] at (2.0,3.75) {$t=3$};
      \node[lec12cu] (b11) at (2.0,3.2) {$v_{11}$};
      \node[lec12tr] (b3) at (1.6,1.9) {$v_3$}; \node[lec12tr] (b1) at (2.4,1.9) {$v_1$};
      \node[lec12pr] (b8) at (1.6,0.6) {$v_8$}; \node[lec12pr] (b7) at (2.4,0.6) {$v_7$};
      \draw[dedge] (b11)--(b3); \draw[dedge] (b11)--(b1);
      \draw[dedge] (b3)--(b8); \draw[dedge] (b1)--(b7);
      \node[font=\small] at (4.6,3.75) {$t=5$};
      \node[lec12cu] (c11) at (4.6,3.2) {$v_{11}$};
      \node[lec12tr] (c1) at (3.8,1.9) {$v_1$}; \node[lec12tr] (c3) at (4.6,1.9) {$v_3$};
      \node[lec12tr] (c5) at (5.4,1.9) {$v_5$};
      \node[lec12pr] (c7) at (3.8,0.6) {$v_7$}; \node[lec12pr] (c8) at (4.6,0.6) {$v_8$};
      \node[lec12pr] (c9) at (5.4,0.6) {$v_9$};
      \draw[dedge] (c11)--(c1); \draw[dedge] (c11)--(c3); \draw[dedge] (c11)--(c5);
      \draw[dedge] (c1)--(c7); \draw[dedge] (c3)--(c8); \draw[dedge] (c5)--(c9);
      \node[lec12cap, anchor=west] at (6.0,3.2) {customer};
      \node[lec12cap, anchor=west] at (6.0,1.9) {transaction};
      \node[lec12cap, anchor=west] at (6.0,0.6) {product};
      \node[lec12cap] at (2.7,-0.85) {(c) computation graphs for different time $t$};
    \end{scope}
  \end{tikzpicture}
  \caption{GNNs on temporal graphs: transactions $v_1$--$v_6$ (blue) link products
  $v_7$--$v_9$ (orange) and customers $v_{10}$--$v_{12}$ (green); the computation graph of
  $v_{11}$ grows with $t$ (slide~52).}
  \label{fig:lec12-temporal}
\end{figure}

\subsection{GNN vs feature engineering}\label{sec:lec12-vsfe}
\slides{53}

\begin{itemize}
  \item \textbf{GNN-based features:} embeddings learned along user $\to$ sales $\to$
    products, \textbf{vs.}
  \item \textbf{Hand-engineered features:}
    \begin{itemize}
      \item \texttt{SUM(TRANSACTIONS.Price)} over (-30, 0) days;
      \item \texttt{AVG(TRANSACTIONS.Price)} over (-30, 0) days; \dots
    \end{itemize}
\end{itemize}
\begin{keyidea}
  GNN aggregation is \textbf{learnable version of hand-crafted features}! GNNs give better
  performance by learning optimal features.
\end{keyidea}

\subsection{SQL joins vs graph edges}\label{sec:lec12-joins}
\slides{54--56}

\begin{itemize}
  \item \textbf{Definitions:} a table $R$ is a set of entities $R=\{r_1,\dots,r_n\}$;
    each entity is a tuple $r_i=(r_{i1},r_{i2},\dots,r_{in})$; given two tables $R,S$, a
    \textbf{join} operation is a \textbf{subset} of a Cartesian product (formula left blank
    on slide~54); aggregation will specify which rows are kept.
  \item \textbf{Example tables:} Employee (LastName, DepartmentID) = (Rafferty, 31),
    (Jones, 33), (Heisenberg, 33), (Robinson, 34), (Smith, 34), (Williams, NULL);
    Department (DepartmentID, DepartmentName) = (31, Sales), (33, Engineering),
    (34, Clerical), (35, Marketing).
  \item \textbf{Ex.\,1: Cross Join:} \texttt{SELECT * FROM employee INNER JOIN department
    ON 1=1;} returns all 24 (employee, department) pairs (Williams with DepartmentID
    NULL included).
  \item \textbf{Ex.\,2: Inner Join:}
\end{itemize}
\begin{equation}\label{eq:lec12-join}
  R\bowtie S\defeq\{r\cup s\mid r\in R\wedge s\in S\wedge \mathrm{Fun}(r,s)\},
  \qquad \mathrm{Fun}(r,s):\ r.pkey=s.fkey
\end{equation}
\begin{itemize}
  \item[] {\small\ttfamily SELECT employee.LastName, employee.DepartmentID,\\
    \hspace*{1em}department.DepartmentName FROM employee INNER JOIN department ON\\
    \hspace*{1em}employee.DepartmentID = department.DepartmentID;}\\
    $\to$ (Robinson, 34, Clerical), (Jones, 33, Engineering), (Smith, 34, Clerical),
    (Heisenberg, 33, Engineering), (Rafferty, 31, Sales); Williams and Marketing drop out.
  \item \textbf{Connection between SQL joins and graph edges:} viewing $R\times S$ as a
    matrix, the kept entries are
\end{itemize}
\begin{equation}\label{eq:lec12-adj}
  a_{ij}=\begin{cases}1, & \text{if } r_i.pkey=s_j.fkey\\ 0, & \text{otherwise}\end{cases}
\end{equation}
\begin{itemize}
  \item[] $\Rightarrow$ \textbf{these are exactly the pkey-fkey edges!}
  \item \textbf{GNNs perform a JOIN+AGG} (slide~56). Input: users 1, 2, 3 and sales
    (USER\_ID, VALUE, DATE) = (1, 20, 01-01), (1, 80, 01-02), (2, 15, 01-01), (3, 20,
    01-02), (3, 80, 01-03); each sale is linked to its user.
  \item \textbf{GNN can learn:} \texttt{SELECT SUM(VALUE) FROM SALES WHERE DATE > 01-01
    GROUP BY USER\_ID} $\to$ (USER\_ID, SUM(VALUE)) = (1, 80), (2, 0), (3, 100), via
\end{itemize}
\begin{equation}\label{eq:lec12-joinagg}
  \bigoplus_{w\in\mathcal{N}(v)}\vt(t_w)\cdot\vx_w
\end{equation}
\begin{itemize}
  \item[] $\bigoplus$: \textbf{learnable aggregation}; $\vt(t_w)$: \textbf{temporal
    embedding}; $\vx_w$: \textbf{fact representation}.
\end{itemize}

\subsection{Benefits of GNNs}\label{sec:lec12-benefits}
\slides{57--61}

Examples on the schema Users, Sessions, Sales, Products, Views:
\begin{itemize}
  \item \textbf{GNNs learn how to aggregate information:}
    \begin{itemize}
      \item they can discard neighboring node information that is irrelevant for the
        given downstream task;
      \item they can detect fine-grained patterns within local neighborhoods (e.g.,
        buying behavior over the last year).
    \end{itemize}
  \item \textbf{GNNs can exchange information \emph{across} training examples:}
    \begin{itemize}
      \item instead of treating examples as isolated, there now exists an
        inter-dependency \emph{between} entities (e.g., users with similar features, users
        with similar behavior);
      \item GNN can \emph{use} these features to enrich an entity's representation (``How
        is the buying behavior of users that have bought similar products?'' / ``\dots\
        that have looked at similar products?'').
    \end{itemize}
  \item \textbf{Multi-hop reasoning \emph{across} table boundaries} can catch information
    which is \textbf{hard} to pre-compute \emph{beforehand} (training labels $\to$ Users;
    ``Which user has bought which product?'', ``How active was the user in the past?'').
  \item \textbf{First-class temporal support} (\cref{fig:lec12-sampling}):
    \begin{itemize}
      \item capture fine-grained relative and seasonal features via \textbf{temporal
        embeddings};
      \item avoid \textbf{data leakage} via \textbf{temporal sampling} directly during
        data loading.
    \end{itemize}
  \item \textbf{Full vision described in paper:} Fey, Hu, Huang, Lenssen, Ranjan,
    Robinson, Ying, You, Leskovec, \emph{Relational Deep Learning: Graph Representation
    Learning on Relational Tables} (Stanford, Kumo.AI, MPI for Informatics, Yale, UIUC).
    Available at \texttt{https://relbench.stanford.edu}.
\end{itemize}

\begin{figure}[htbp]
  \centering
  \begin{tikzpicture}[scale=0.85, transform shape]
    \node[lec12smp] (f1) at (-1.6,1.6) {}; \node[lec12smp] (f2) at (0.8,1.8) {};
    \node[lec12smp] (f3) at (-0.3,1.1) {}; \node[lec12smp] (f4) at (-1.5,0.2) {};
    \node[lec12smp] (f5) at (-0.9,-0.9) {};
    \node[lec12smp, fill=Purple!35] (tp) at (0,0) {};
    \node[font=\small\bfseries] at (0.25,-0.55) {T+1};
    \draw[densely dashed, black!70] (f1)--(f3) (f2)--(f3) (f3)--(tp) (f4)--(tp) (f5)--(tp);
    \node[lec12smp, fill=Purple!35] (t) at (2.6,0.75) {};
    \node[font=\small\bfseries] at (2.6,0.2) {T};
    \draw[densely dashed, thick, black!70] (tp)--(t);
    \node[lec12smp] (a1) at (4.0,1.4) {}; \node[lec12smp] (a2) at (5.3,1.6) {};
    \node[lec12smp] (b1) at (4.4,-0.1) {}; \node[lec12smp] (b2) at (5.4,0.6) {};
    \node[lec12smp] (b3) at (5.5,-0.8) {};
    \node[font=\small\bfseries] at (4.4,-0.7) {T-1};
    \draw[thick, Purple!70, -{Stealth[length=1.6mm]}] (a1)--(t);
    \draw[thick, Purple!70, -{Stealth[length=1.6mm]}] (a2)--(a1);
    \draw[thick, Purple!70, -{Stealth[length=1.6mm]}] (b1)--(t);
    \draw[thick, Purple!70, -{Stealth[length=1.6mm]}] (b2)--(b1);
    \draw[thick, Purple!70, -{Stealth[length=1.6mm]}] (b3)--(b1);
    \node[font=\small, BrickRed, align=center] at (-3.2,0.4) {Not available\\for sampling};
    \node[font=\small, BrickRed, align=center] at (7.1,0.4) {Available for\\sampling};
  \end{tikzpicture}
  \caption{Temporal sampling for a seed node at time T: neighbors from the past (T-1) can be
  sampled, the future (T+1) cannot (slide~60).}
  \label{fig:lec12-sampling}
\end{figure}

% ---------------------------------------------------------------------
\section{RelBench}\label{sec:lec12-relbench}

\subsection{Enabling research on RDL}\label{sec:lec12-enabling}
\slides{63--64}

\begin{itemize}
  \item \term{RelBench} (Relational Deep Learning Benchmark) is more than just a
    collection of databases: RelBench Relational Databases $\to$ RelBench Data, Task,
    Graph Loaders $\to$ Your Relational DL Model (PyG, PyTorch Frame) $\to$ RelBench
    Evaluator $\to$ RelBench Leaderboard.
    \begin{itemize}
      \item automatically download datasets; load database and task tables;
      \item standardized evaluation protocol: \textbf{prevents temporal leakage from test
        set};
      \item framework-agnostic data structures: use your favorite ML stack!
    \end{itemize}
  \item \textbf{PyF and PyG integration:} load as a PyG graph; train GNN end-to-end;
    temporal neighbor sampling; use PyTorch Frame to encode tables.
\end{itemize}

\subsection{Datasets and tasks}\label{sec:lec12-datasets}
\slides{65--68}

\begin{itemize}
  \item \textbf{7 diverse datasets:} E-Commerce -- rel-amazon, rel-avito, rel-hm; Social
    -- rel-event, rel-stack; Sports -- rel-f1; Medical -- rel-trial.
  \item \textbf{Rich schemas:} \textbf{3} to \textbf{15} tables; \textbf{74k} to
    \textbf{41M} rows in a DB; \textbf{15} to \textbf{140} columns in a DB; time span
    from \textbf{2 weeks} to \textbf{55 years}. Examples: rel-amazon -- review(review\_text,
    summary, review\_time, rating, verified, customer\_id, product\_id),
    customer(customer\_id, customer\_name), product(product\_id, brand, title,
    description, price, category); rel-trial (slide~36).
  \item \textbf{30 real-world predictive tasks} (\cref{tab:lec12-tasks}).
  \item \textbf{See website to get started:} \texttt{https://relbench.stanford.edu}
    (``Open benchmark for machine learning over relational databases''; realistic
    databases, flexible data loaders, evaluators; beta testing phase).
\end{itemize}

\begin{table}[htbp]
  \centering\small
  \caption{RelBench tasks shown on slide~67 (``\dots'' = more tasks).}
  \label{tab:lec12-tasks}
  \begin{tabular}{@{}lll@{}}
    \toprule
    Entity classification & Entity regression & Recommendation\\
    \midrule
    rel-amazon: user-churn, item-churn & rel-amazon: user-ltv, item-ltv
      & rel-amazon: user-item-purchase\\
    rel-stack: user-badge & rel-avito: ad-ctr & rel-avito: user-ad-visit\\
    rel-trial: study-outcome & rel-f1: driver-position
      & rel-stack: user-post-comment,\\
    \dots & \dots & \quad post-post-related; \dots\\
    \bottomrule
  \end{tabular}
\end{table}

% ---------------------------------------------------------------------
\section{GNN vs expert data scientist}\label{sec:lec12-expert}

\subsection{Setup}\label{sec:lec12-setup}
\slides{70--72}

\begin{itemize}
  \item \textbf{Expert data scientist} (Alejandro): recruited experienced data scientist;
    5 years in industry (specializing in financial databases); responsible for full model
    building lifecycle.
  \item \textbf{Representative RelBench task}: \textbf{Stack Exchange database} (tables postLinks, posts, badges, votes, postHistory, users,
    comments). \textbf{Q: Will a user be active in the next 6 months?}
\end{itemize}

\subsection{Expert data scientist workflow}\label{sec:lec12-workflow}
\slides{73--81}

\begin{table}[htbp]
  \centering\small
  \caption{Expert data scientist workflow, manual work per stage (slides~73--79).}
  \label{tab:lec12-workflow}
  \begin{tabularx}{\linewidth}{@{}>{\raggedright\arraybackslash}p{3.7cm}r>{\raggedright\arraybackslash}X@{}}
    \toprule
    Stage & Manual work & Details on the slides\\
    \midrule
    Exploratory Data Analysis (EDA) & 4hrs & example observations: activity is seasonal
      (new posts over time); comments follow power law (dist.\ of num comments per
      user); negative votes are infrequent (dist.\ of votes by post)\\
    Feature Ideation & 0.5hr & example features: weeks\_since\_last\_comment,
      badge\_score, location\_is\_null\\
    SQL query writing & 5hr & 100s of lines of code (\texttt{create table
      churn\_feats\_train as \dots}, joining labels with badges on
      \texttt{labels.timestamp > badges.Date})\\
    XGBoost hparam sweep & 2hr & \\
    SHAP (feature importance analysis) & 1hr & \\
    \bottomrule
  \end{tabularx}
\end{table}

\begin{itemize}
  \item \textbf{Manual work:} EDA $\to$ feature ideation $\to$ SQL query writing $\to$
    XGBoost hparam sweep $\to$ SHAP, with times in \cref{tab:lec12-workflow}.
  \item \textbf{SHAP feature importance} (top of the plot: weeks\_since\_last\_comment,
    badge\_score, last\_q\_weeks\_ago, num\_posts\_commented, last\_a\_weeks\_ago,
    num\_badges, location\_is\_null, about\_me\_length, \dots). Selected observations:
    \begin{itemize}
      \item website ``seniority'' \textbf{predictive} (total number of comments, badge
        score etc.);
      \item time since last active / commented \textbf{is predictive};
      \item completed bio (about me, location etc.) \textbf{is predictive};
      \item number of positive/negative votes \textbf{not predictive};
      \item interacting with ``senior'' community members \textbf{not predictive}.
    \end{itemize}
  \item \textbf{Final list of features:} 54 features with Label Corr., Label MI, NaN \%
    (highest label correlation: num\_badges 0.229, num\_questions\_last\_6mo 0.209,
    badge\_score 0.206; most negative: last\_a\_weeks\_ago $-0.394$,
    weeks\_since\_last\_comment $-0.392$, last\_q\_weeks\_ago $-0.346$).
    \textbf{$\sim$12 hours of high-quality expert work.}
\end{itemize}

% ---------------------------------------------------------------------
\section{Results}\label{sec:lec12-results}

\subsection{Stack Exchange: will a user be active in the next 6 months?}\label{sec:lec12-stack}
\slides{83--86}

\begin{itemize}
  \item \textbf{Naive baseline:} XGBoost (no feature engineering) on the users table (Id,
    AccountId, DisplayName, Location, WebsiteUrl, AboutMe, CreationDate).
  \item Test ROCAUC and effort in \cref{tab:lec12-results}. Expert data scientist:
    \textbf{12 hours manual work, $\sim$682 lines of code, 1hr model training (2 days to
    model)}; GNN: \textbf{$<$1 hour manual work, 54 lines of code, 1hr model training}
    (*work measured as the marginal effort to solve a new task).
  \item Performance comparison (slide~86): AUCROC 0.90 (GNN) vs 0.86 (expert); hours
    human work 1.00 vs 11.00; lines of code 54.00 vs 682.00; hours model training 1.00 vs
    1.00. (The chart's 11.00 hours differs from ``$\sim$12 hours'' on slide~81 and ``12
    hours'' on slides~84--85.)
\end{itemize}

\begin{table}[htbp]
  \centering\small
  \caption{GNN vs expert data scientist. Top: Stack Exchange ``will a user be active in
  the next 6 months?'', Test ROCAUC and effort (slides~83--85); bottom: AUCROC on more
  RelBench tasks (slide~87).}
  \label{tab:lec12-results}
  \begin{tabular}{@{}lccc@{}}
    \toprule
    Task & XGBoost (no feature eng.) & Expert Data Scientist & GNN\\
    \midrule
    User active in next 6 months & 0.6842 & 0.8643 & \textbf{0.8945}\\
    \quad manual work & -- & 12 hours & $<$1 hour\\
    \quad lines of code & -- & $\sim$682 & 54\\
    \quad model training & -- & 1hr & 1hr\\
    \midrule
    stackex-engage (AUCROC) & -- & 86.43 & \textbf{89.50}\\
    stackex-badges (AUCROC) & -- & 86.14 & \textbf{88.19}\\
    amazon-churn (AUCROC) & -- & 67.55 & \textbf{70.46}\\
    amazon-product-churn (AUCROC) & -- & 81.67 & \textbf{82.98}\\
    \bottomrule
  \end{tabular}
\end{table}

\subsection{More RelBench tasks}\label{sec:lec12-more}
\slides{87}

\begin{itemize}
  \item Performance comparison on stackex-engage, stackex-badges, amazon-churn and
    amazon-product-churn (\cref{tab:lec12-results}, bottom): \textbf{GNN beats the expert
    data scientist on all four tasks}.
\end{itemize}

% ---------------------------------------------------------------------
\bigskip\phantomsection
\addcontentsline{toc}{section}{Lecture 12 at a glance}
\begin{summarybox}[Lecture 12 at a glance]
  \begin{itemize}
    \item Data is stored in \textbf{relational databases}; classic ML flattens them
      by manual feature engineering (6-12 months; arbitrary features, limited data,
      leakage). LLMs: not relational, not sequence modeling, not trained to predict the
      future, limited context.
    \item \textbf{RDL:} end-to-end deep learning on relational tables, no feature
      engineering; predictive tasks become graph representation learning. Inheritor of
      statistical relational learning; unlike tabular ML, uses relations between tables.
    \item Database $\calT=\{T_1,\dots,T_n\}$, links $\mathcal{L}\subseteq\calT\times\calT$;
      \textbf{schema graph} = high-level structure; \textbf{relational entity graph}:
      nodes $\calV=\bigcup_{T\in\calT}T$ (all rows), edges = matching primary/foreign
      keys; entities have features.
    \item \textbf{Task} = entity + label + time, given by a \textbf{training table} (entity
      ID, time, labels); only use information up to the time of label.
    \item \textbf{Relational GNN:} time-dependent computation graphs, message+aggregation
      and sampling; aggregation = learnable hand-crafted features; inner join
      $a_{ij}=1$ iff $r_i.pkey=s_j.fkey$ = pkey-fkey edges; JOIN+AGG
      $\bigoplus_{w\in\mathcal{N}(v)}\vt(t_w)\cdot\vx_w$.
    \item \textbf{RelBench:} 7 datasets, 30 tasks, loaders + evaluator without temporal
      leakage; PyG + PyTorch Frame.
    \item \textbf{GNN vs expert} (Stack Exchange): 0.8945 vs 0.8643 (XGBoost baseline
      0.6842) with $<$1 h vs 12 h manual work and 54 vs $\sim$682 lines of code; GNN also
      wins on stackex-engage, stackex-badges, amazon-churn, amazon-product-churn.
  \end{itemize}
\end{summarybox}
